% !TeX root = ../../../main.tex
\subsection{Survival Modeling with Language Model Embeddings}
\label{ch07:apd:encoder-survival}


% !TeX root = ../../../main.tex
\begin{table}[htbp]
\caption{Hyperparameter grid search ranges for our survival models.}\ChSevenFloatLabel{ch07:tab:hyperparams}
\ChSevenTableBodyStart
\begin{ChSevenFloatBody}

\centering
\begin{adjustbox}{max width=\columnwidth}
\begin{tabular}{ll}
\toprule
\textbf{Model} & \textbf{Hyperparameters} \\
\midrule
RSF & \texttt{n\_estimators} $\in$ \{20,50,100,200,500,1000\} \\
    & \texttt{min\_samples\_split} $\in$ \{2,10,20,50,100\} \\
    & \texttt{min\_samples\_leaf} $\in$ \{1,5,10,20\} \\
DeepSurv/DeepHit & \texttt{num\_nodes} $\in$ \{32,64,128,256,512,768,1024,1200,\\
& 1536,1800,2048,3072,4096\} \\
    & \texttt{dropout} $\in$ \{0.1,0.2,0.4,0.5\} \\
\bottomrule
\end{tabular}
\end{adjustbox}



\end{ChSevenFloatBody}
\end{table}


To evaluate the prognostic value of textual information encoded in large language models (LLMs), we adopt a two-stage framework: (1) extraction of fixed-dimensional sequence embeddings from various pre-trained LLMs, and (2) downstream survival modeling using these embeddings as covariates.

In the first stage, we process each textual time series (which is converted to a textual context of \emph{"(time) clinical event [SEP] ... [SEP] (time) clinical event [SEP]"}) using a suite of LLMs to obtain dense vector representations that summarize the content of the input sequence. 
For models belonging to the encoder family, including \textit{bert-base-uncased}, \textit{roberta-base}, \textit{deberta-v3-small}, \textit{ModernBERT-base}, and \textit{ModernBERT-large}, we extract the final hidden state corresponding to the \textit{[CLS]} token, which is conventionally used to represent the entire sequence in classification tasks. 
This token-specific embedding serves as a compact, sequence-level representation.

For decoder-based models that do not utilize a dedicated \texttt{[CLS]} token, including \textit{DeepSeek-R1-Distill-Llama-70B}, \textit{Llama-3.3-70B-Instruct}, \textit{DeepSeek-R1-Distill-Llama-8B}, and \textit{Llama-3.1-8B-Instruct}, we compute the mean-pooled embedding over the last hidden states of all non-padding tokens. 
This pooling strategy yields a fixed-length vector that captures the overall semantic content of the input while mitigating the impact of padding artifacts.

In the second stage, these embedding vectors are used as input covariates to three survival models, each designed to capture time-to-event dynamics in different ways:

\begin{itemize}
    \item \textbf{Random Survival Forest (RSF):} 
    A nonparametric ensemble method based on decision trees, capable of handling complex interactions and high-dimensional inputs. 
    The model is tuned over the number of trees (\texttt{n\_estimators}), the minimum number of samples required to split an internal node (\texttt{min\_samples\_split}), and the minimum number of samples at a leaf node (\texttt{min\_samples\_leaf}); see Table~\ref{ch07:tab:hyperparams} for search ranges.

    \item \textbf{DeepSurv:} 
    A neural-network generalization of the Cox proportional hazards model. 
    The architecture comprises fully connected layers with varying widths (\texttt{num\_nodes}) and dropout rates; see Table~\ref{ch07:tab:hyperparams} for search ranges.
The number of training epochs is fixed.

    \item \textbf{DeepHit:} 
    A multi-task neural network that jointly models the discrete hazard and survival functions. 
    As with DeepSurv, we perform a grid search over hidden layer sizes and dropout rates while keeping the number of epochs fixed; see Table~\ref{ch07:tab:hyperparams}.
\end{itemize}

Hyperparameter tuning for each survival model is performed using a validation split from the training set. 
For each combination of language model and survival model, we conduct grid search over the respective hyperparameter space, selecting the configuration that maximizes the time-dependent concordance index \citep{ch07:antolini2005time} on the validation set. 
Final evaluation is conducted on held-out test sets (\textit{T2S2\_test}, \textit{sepsis-10}, \textit{sepsis-100}), and results are aggregated across survival metrics to assess model performance.

This framework allows for a systematic evaluation of how different pre-trained LLMs, and their associated embedding strategies, contribute to downstream survival prediction tasks.
%%%%%%%%%%%%%%%%%%%%%%%%%%%%%%%%%%%%%%%%%%%%%%%%
%%%%%%%%%Appendix F: Forecasting tables
%%%%%%%%%%%%%%%%%%%%%%%%%%%%%%%%%%%%%%%%%%%%%%%
\FloatBarrier

% \section{Forecasting performance on \textit{DeepSeek-R1} and \textit{Llama-3.3-70B} textual time series annotations}
% \label{apd:full_tables}

